\begin{exo}{}
	On considère un parallélogramme $ABCD$ de centre $O$ et le triangle équilatéral $BCE$.


	$G$, $H$ et $I$ sont les milieux respectifs de $[BE]$, de $[CE]$ et de $[BC]$.

	\begin{enumerate}
		\item Faire une figure à main levée.
		\item Indiquer trois couples de vecteurs égaux deux à deux en justifiant à l'aide de propriétés vectorielles du cours.
	\end{enumerate}
\end{exo}
